% Clasificación exacta del entrelazador horizonte--TFD--Perron.

\section{Clasificación espectral del entrelazador entre horizonte, purificación termocampo doble y dinámica de Perron}
\label{sec:ivv-horizon-tfd-perron-classification}

Sean \(\mathcal H_{\mathrm H}\), \(\mathcal H_{\mathrm{TFD}}\) y
\(\mathcal H_{\mathrm P}\) los subespacios cíclicos generados,
respectivamente, por los microestados de horizonte, la purificación
termocampo doble y el operador positivo de Perron. Denotemos por
\(A_{\mathrm H}\), \(A_{\mathrm{TFD}}\) y \(A_{\mathrm P}\) sus operadores
autoadjuntos restringidos y por \(m_X(\lambda)\) la multiplicidad espectral
de \(\lambda\) en el sector \(X\).

\begin{theorem}[Existencia y unicidad por calibres del entrelazador espectral]
\label{thm:ivv-horizon-tfd-perron-classification}
Existe una isometría sobreyectiva
\[
 U:\mathcal H_{\mathrm H}\longrightarrow\mathcal H_{\mathrm{TFD}}
 \longrightarrow\mathcal H_{\mathrm P}
\]
que entrelaza los \(3\) operadores si y sólo si, después de aplicar las
normalizaciones declaradas de energía y memoria,
\[
 \operatorname{spec}(A_{\mathrm H})
 =\operatorname{spec}(A_{\mathrm{TFD}})
 =\operatorname{spec}(A_{\mathrm P}),
 \qquad
 m_{\mathrm H}(\lambda)=m_{\mathrm{TFD}}(\lambda)=m_{\mathrm P}(\lambda)
\]
para todo \(\lambda\) del espectro común. Cuando estas condiciones se
cumplen, el entrelazador queda determinado salvo una transformación unitaria
en cada subespacio propio de multiplicidad mayor que \(1\).
\end{theorem}

\begin{proof}
Si \(U A_X=A_YU\), la equivalencia unitaria preserva espectro y
multiplicidades. Recíprocamente, para cada \(\lambda\) se eligen bases
ortonormales de los subespacios propios de igual dimensión y se define
\(U_\lambda\) por correspondencia de bases. La suma ortogonal
\(U=\bigoplus_\lambda U_\lambda\) es unitaria e intercambia los operadores.
La libertad residual es exactamente el producto de los grupos unitarios de
los subespacios degenerados.
\end{proof}

El problema de existencia queda así cerrado por clasificación: coincidencia
de espectros y multiplicidades produce el entrelazador; su incumplimiento
demuestra la obstrucción. Una igualdad escalar aislada no sustituye estas
condiciones.

